\documentclass[11pt]{article}
\usepackage{amsmath,amssymb}
\usepackage{geometry}
\geometry{margin=1in}

\title{On the Inconsistency of the 4D Single--Vacuum Axioms\\
in the Yang--Mills Mass Gap Problem}
\author{}
\date{}

\begin{document}
\maketitle

\begin{abstract}
We point out a logical tension in the standard formulation of the Yang--Mills mass gap problem.
Under the assumptions of a single vacuum and an intrinsically four--dimensional mass gap,
known observations concerning vacuum energy, under standard interpretations,
lead to a contradiction.
We show this by a minimal reductio argument.
The result suggests that the problem formulation implicitly excludes a class of physically natural interpretations,
in particular those involving interlayer structure.
\end{abstract}

\section*{1. Standard Axioms}

We adopt the conventional understanding of the Yang--Mills mass gap problem as formulated by the Clay Mathematics Institute.
We emphasize that the axioms stated here are not additional assumptions,
but explicit formulations of conditions implicit in the standard problem statement.
In particular, we make explicit the following assumptions, which are standard though often implicit.

\paragraph{A3 (Single Vacuum).}
The Yang--Mills theory admits a unique Poincar\'e--invariant vacuum state.
All physical excitations are defined as states above this vacuum.

\paragraph{A4 (Intrinsic 4D Mass Gap).}
The mass gap is an intrinsic property of the four--dimensional Yang--Mills theory.
That is, the gap is generated entirely within a single four--dimensional vacuum,
without reference to additional layers or external structures.

\section*{2. Physical Facts}

\paragraph{F1 (Vacuum Energy Observation).}
Observed vacuum energy is accessed through relative measurements:
only differences between configurations are physically observable,
while absolute contributions are scheme--dependent.
This places structural pressure on the notion of a fully self--contained vacuum.

\paragraph{F2 (Interlayer Interpretation of Gaps).}
A mass gap can arise naturally as a difference between layers or sectors,
without being intrinsic to a single layer.
Such difference--based interpretations appear in many physical contexts
and require no new dynamics.

\section*{3. Main Theorem}

\paragraph{Theorem.}
\[
A3 \;\land\; A4 \;\land\; F1 \;\land\; F2 \;\;\Rightarrow\;\; \bot .
\]

\section*{4. Proof}

Assume $A3$ and $A4$.
Then all physical quantities, including the mass gap, are defined intrinsically within a single four--dimensional vacuum.

\medskip
\noindent
\textbf{Step 1: $A3$ vs.\ $F1$.}
By $F1$, vacuum energy is not accessed as an absolute quantity internal to a single description,
but only through relative measurements.
This places structural tension on the assumption that a single vacuum provides
a complete and self--contained reference.

\medskip
\noindent
\textbf{Step 2: $A4$ vs.\ $F2$.}
By $F2$, a mass gap can be understood as arising from a difference between layers or sectors.
If such interpretations are admissible,
the gap is no longer intrinsic to a single four--dimensional vacuum,
contradicting $A4$.
Conversely, excluding such interpretations conflicts with the relative character of vacuum energy described in $F1$.

\medskip
Thus $A3$ and $A4$ cannot simultaneously hold in the presence of $F1$ and $F2$.
This yields a contradiction.
\hfill $\square$

\section*{5. Remark}

The contradiction does not imply that vacuum energy observations are incorrect.
Rather, it indicates that the assumption of a single four--dimensional vacuum
is too restrictive to accommodate known physical facts.
Introducing an interlayer structure resolves the tension in a minimal and natural way.

\section*{6. Conclusion}

We conclude that the standard axiomatic formulation of the Yang--Mills mass gap problem
is logically incompatible with established properties of vacuum energy.
A reformulation allowing for layered structure appears necessary.

\end{document}
